\section{Experiments}
\label{sec:experiments}
% DRAFT (Claude): awaits the author's rewrite. Results stay as TODOs until the
% pilot and judge runs are in; settings match configs/pilot.yaml.

\subsection{Data}
Trajectories come from the \texttt{random} split of LeanDojo Benchmark 4
\citep{yang2023leandojo}, which records the goal state before and after
every tactic of each traced Mathlib theorem. Most traced proofs are one to
four tactics long, too short for an area change or for a
Chapman--Kolmogorov test at lags of a few steps, so we keep proofs with at
least 8 steps and sample up to 20{,}000 of them. \todo{number of proofs and
steps after filtering; share of Mathlib folders covered.} Lean Workbook
\citep{ying2024leanworkbook} records at most one tactic per problem, so it
contributes goal states for judge validation but no transitions. miniF2F and
PutnamBench contain statements without proofs and serve only as search
targets. As a sanity check, a synthetic set with planted areas runs
through the same pipeline end to end.

\begin{table}[t]
  \caption{Data. Proof counts are after filtering to at least 8 steps;
  benchmark sets contribute statements only. \todo{fill from the pilot run.}}
  \label{tab:data}
  \centering\small
  \setlength{\tabcolsep}{5pt}
  \resultstable{table_data}
\end{table}

\subsection{Is the model Markovian?}

\begin{figure}[t]
  \centering
  \maybefigure{implied_timescales}{0.47\linewidth}\hfill
  \maybefigure{transition_matrix}{0.47\linewidth}
  \caption{Left: implied timescales against the lag $\lag$, both in proof
  steps; bands are bootstrap intervals over proofs, and the gray region marks
  processes faster than the lag. Right: transition matrix between
  macrostates at the chosen lag. \todo{lag and number of macrostates.}}
  \label{fig:its}
\end{figure}

\begin{figure}[t]
  \centering
  \maybefigure{chapman_kolmogorov}{0.9\linewidth}
  \caption{Chapman--Kolmogorov test. Lines show the model's prediction
  $\T(\lag)^k$; points show $\T(k\lag)$ estimated directly from the data,
  with bootstrap error bars.}
  \label{fig:ck}
\end{figure}

\todo{which lag the timescales converge at; how many slow processes are
resolved; whether the CK test passes, and for which macrostate pairs it
fails.}

\subsection{Do macrostates match areas?}

\begin{figure}[t]
  \centering
  \maybefigure{judge_confusion}{0.47\linewidth}\hfill
  \maybefigure{area_flow}{0.47\linewidth}
  \caption{Left: macrostates (rows) against the judges' consensus MSC code
  (columns), row-normalized; NMI $=$ \resMacroNMI. Right: the \resFlowsShown
  largest flows between areas within proofs, covering \resFlowCoverage\% of
  area changes.}
  \label{fig:areas}
\end{figure}

We call the area hypothesis supported if the macrostates reach an NMI of at
least 0.15 against the judges' consensus labels and beat the Mathlib-folder
baseline. Both conditions were fixed before any labels were collected.
\todo{judge agreement ($\kappa$) and accuracy on known codes; NMI and ARI
for $m = 10$ and the sweep over $m$; next-area prediction against the
``stay'' and marginal baselines.}

\begin{table}[t]
  \caption{Agreement with the judges' consensus MSC codes. Brackets are
  95\% bootstrap intervals over proofs; AMI corrects NMI for the number of
  groups. The rows above the second rule are scored on the same labelled
  steps. The folder and namespace rows label each step by its Mathlib folder
  or the first components of its theorem name; the $k$-means rows cluster the
  same goal-state features, or the same microstate centers, into $m$ groups
  with no dynamics; the judges' proof label gives every step the judges'
  own whole-proof label, an oracle; the random partition keeps the
  macrostate sizes. The null rows give the macrostates' mean score when
  their labels are shuffled over all steps, swapped between proofs of the
  same length, or shuffled within each proof, with $p = $ \resNullRowP,
  \resNullProofP{} and \resNullWithinP{} for the observed NMI. The last row scores the synthetic set against its planted
  areas. \todo{fill from the judge run.}}
  \label{tab:areas}
  \centering\small
  \resultstable{table_areas}
\end{table}
